%% kapa: exam=1 round=27 year=2016 period=1 form=A subject=real_estate subject_ko=부동산학원론 seq=4 number=84
%% source: Q-Net 감정평가사 기출문제 (cst00309) artlSeq=5169402 file=extracted/1008/5169402/제27회 감정평가사 1차 문제/제27회 감정평가사 1교시문제(A형).pdf
%% conversion: kordoc 4.24.51 -> tools/gen1.py (auto); review=proofread-sonnet; flags=none
\begin{question}{84}
부동산의 특성에 관한 설명으로 옳지 않은 것은?
\begin{choices}
  \item 부동성으로 인해 부동산활동을 국지화시키고 임장활동을 배제한다.
  \item 토지는 물리적인 측면에서는 영속성을 가지나, 경제적 가치는 주변상황의 변화에 의하여 하락될 수 있다.
  \item 영속성으로 인해 토지는 감가상각에서 배제되는 자산이다.
  \item 개별성으로 인해 부동산활동이 구체적이고 개별적으로 전개되며, 부동산시장에서 정보의 중요성이 증대된다.
  \item 용도의 다양성으로 인해 토지이용결정과정에서 용도가 경합할 경우, 최유효이용할 수 있는 방안을 도출하여 실행하게 한다.
\end{choices}
\end{question}
\answer{1}
